\documentclass{article}
\usepackage[UTF8]{ctex}
\usepackage{amsmath,amssymb,amsfonts}
\usepackage{enumitem}
\usepackage{geometry}
\geometry{a4paper, margin=1in}
% === 错题本排版优化（不参与题目标识） ===
\setlength{\parindent}{0pt}
\setlength{\parskip}{0.35em plus 0.10em minus 0.08em}
\setlength{\emergencystretch}{3em}
\setlist{itemsep=0.25em, topsep=0.40em, parsep=0pt}
\allowdisplaybreaks[2]
% === 错题本排版优化结束 ===
\begin{document}
\section*{华东理工大学2026年高等代数真题}
\section*{47．华东理工大学2026年高等代数试卷}

\subsection{题目1}
一．求所有整数 \(m\) ，使多项式 \(f(x)=x^{5}+m x+1\) 在有理数域上可约．

\subsection{题目2}
二．设 \(n(n>2)\) 阶行列式 \(D\) 所有元素为1或－1，求证：\(D\) 的绝对值 \(|D| \leq(n-1)!(n-1)\) 。

\subsection{题目3}
三。设 \(a, \beta\) 为不同的 \(n\) 维实列向量，若 \(\|\alpha\|=\|\beta\|\) ，证明：存在 \(n\) 阶方阵 \(H=I_{n}-2 u u^{\prime}\) ，使得\(H \alpha=\beta\) ，其中 \(u\) 为某个单位向量。

\subsection{题目4}
四．设四元齐次线性方程组\\
（1）：\(\left\{\begin{array}{l}x_{1}+x_{2}=0, \\ x_{3}-x_{4}=0 .\end{array}\right.\)\\
又已知某齐次线性方程组（2）的基础解系为

\[
\eta_{1}=\left(\begin{array}{l}
0 \\
1 \\
1 \\
0
\end{array}\right), \eta_{2}=\left(\begin{array}{c}
-1 \\
2 \\
2 \\
1
\end{array}\right)
\]

求方程组（1）与（2）的公共解。

\subsection{题目5}
五．设

\[
\begin{aligned}
& \alpha_{1}=(1,1,0,0)^{\prime}, \alpha_{2}=(0,1,1,0)^{\prime}, \alpha_{3}=(0,0,1,1)^{\prime} \\
& \beta_{1}=(1,0,1,0)^{\prime}, \beta_{2}=(0,2,1,1)^{\prime}, \beta_{3}=(1,2,1,2)^{\prime}
\end{aligned}
\]

求向量组 \(\alpha_{1}, \alpha_{2}, \alpha_{3}\) 和向量组 \(\beta_{1}, \beta_{2}, \beta_{3}\) 所生成的两个线性空间的和与交的维数与一组基。

\subsection{题目6}
六．设 \(M\) 是数域 \(F\) 上的 \(n\) 阶方阵，\(f(x), g(x)\) 是数域 \(F\) 上的多项式，\(A=f(M), B=\) \(g(M), W, W_{1}, W_{2}\) 分别为方程组 \(A B X=0, A X=0, B X=0\) 的解空间．若 \(f(x), g(x)\) 互素，求证：\(W=W_{1} \oplus W_{2}\) 。

\subsection{题目7}
七．设

\[
A=\left(\begin{array}{ccc}
1 & a & b \\
0 & \omega & c \\
0 & 0 & \omega^{2}
\end{array}\right)
\]

其中 \(a, b, c\) 为任意数，\(\omega=\frac{-1+\sqrt{-3}}{2}\) ，求 \(A^{100}\) 及 \(A^{-1}\) ．

\subsection{题目8}
八．设 \(\alpha_{1}, \alpha_{2}, \cdots, \alpha_{m}\) 与 \(\beta_{1}, \beta_{2}, \cdots, \beta_{m}\) 为 \(n\) 维欧氏空间 \(V\) 的两组向量．证明：若

\[
\left(\alpha_{i}, \alpha_{j}\right)=\left(\beta_{i}, \beta_{j}\right), i, j=1,2, \cdots, m .
\]

则由向量组生成的子空间 \(L\left(\alpha_{1}, \alpha_{2}, \cdots, \alpha_{m}\right)\) 与 \(L\left(\beta_{1}, \beta_{2}, \cdots, \beta_{m}\right)\) 同构。

\subsection{题目9}
九。设 \(A\) 为 \(n\) 阶正定对称矩阵，\(n\) 维实列向量组 \(\alpha, \beta\) 满足 \(\alpha^{\prime} \beta>0\) ，求证：

\[
H=A-\frac{A \beta \beta^{\prime} A}{\beta^{\prime} A \beta}+\frac{\alpha \alpha^{\prime}}{\alpha^{\prime} \beta}
\]

是正定矩阵。

\subsection{题目10}
十．判定实二次型 \(f\left(x_{1}, x_{2}, x_{3}\right)=5 x_{1}^{2}+x_{2}^{2}+5 x_{3}^{2}+4 x_{1} x_{2}-8 x_{1} x_{3}-4 x_{2} x_{3}\) 是否正定．
\end{document}
